%%^^A%%  fontspec-doc-xetex.tex -- part of FONTSPEC <latex3.github.io/fontspec>

\documentclass[a4paper]{l3doc}
\usepackage{fontspec-doc-style}
\showexamplesfalse
\begin{document}

\part{\XeTeX\ 下的字体与特性}
\label{sec:xetex-features}

\section{仅 \XeTeX\ 支持的字体特性}

本节所述特性可用于 \pkg{fontspec} 选择的任意字体。

\subsection{映射}
\label{sec:mapping}

\feat{Mapping} 特性启用 \XeTeX\ 文本映射方案，示例见 \exref{mapping}。

\begin{Xexample}{mapping}{\XeTeX\ 的 \feat{Mapping} 特性。}
  \fontspec{texgyrepagella-regular.otf}[Mapping=tex-text]
  ``!`A small amount of---text!''
\end{Xexample}

同一时刻只能激活一个映射；第二次调用 \feat{Mapping} 将覆盖第一次。
使用 |tex-text| 映射等价于书写 |Ligatures=TeX|；为更好地兼容
\LuaTeX\ 文档，推荐使用后者写法。


\subsection{不同的字体技术：\AAT、OpenType 与 Graphite}\label{sec:renderer-xetex}

\textbf{请注意：自 2020 年起，\XeTeX\ 在 \MacOSX\ 上似乎已无法支持 \AAT\ 字体。}

\XeTeX\ 支持三种排版渲染技术，通过 \feat{Renderer} 字体特性选择。
第一种 \opt{AAT} 仅由 \MacOSX\ 提供。第二种 \opt{OpenType} 是开源 OpenType
解释器；相比 \opt{AAT}，它对 OpenType 特性提供更强支持，尤其是上下文排版。
第三种是 \opt{Graphite}，它是 OpenType 的一种替代方案，针对较少见语言提供
特定特性，并具备更强大的字体选项能力。\opt{OpenType} 的特性已在
\vref{sec:opentype-features} 中讨论；\opt{Graphite} 与 \opt{AAT} 的特性
分别在 \vref{sec:graphite-features} 与 \vref{sec:aat-features} 中讨论。

除非有特殊需要，\feat{Renderer} 特性很少需要显式指定：OpenType 字体会自动使用
\opt{OpenType} 渲染器，\AAT\ 字体默认选择 \opt{AAT}。但有些字体包含多种渲染
技术的字体表（例如随 \MacOSX\ 分发的日文 Hiragino 字体），此时可能更倾向选用
其中之一。

在仅能通过特定渲染器使用的其他字体特性中，\opt{OpenType} 提供 \feat{Script}
与 \feat{Language} 特性，可针对不同字母与语言呈现不同字体行为；这些特性的
描述见 \vref{sec:ot}。{\em 由于这些字体特性可能改变字体实例可选哪些特性，
\pkg{fontspec} 会在所有其他特性之前先选择它们，并自动且不加提示地选用
\opt{OpenType} 渲染器。}


\subsection{竖排}

\XeTeX\ 通过在字体排版时旋转单个字形的能力来支持竖排：
\begin{Verbatim}
  \def\verttext{共産主義者は}
  \fontspec{Hiragino Mincho Pro}
  \verttext

  \fontspec{Hiragino Mincho Pro}[Renderer=AAT,Vertical=RotatedGlyphs]
  \rotatebox{-90}{\verttext}% requires the graphicx package
\end{Verbatim}

并未提供自上而下文字的实际排版支持；有关实现方式的示例，
见 \XeTeX\ 文档中提供的竖排中文示例。


\section{Graphite 渲染器}
\label{sec:graphite-features}

由于 Graphite 渲染器面向较少见的文字与语言设计，这些文字通常具有特定或
独特的需求，Graphite 特性在各字体间并不统一。

目前 \pkg{fontspec} 尚未提供选择 Graphite 字体特性的便捷接口，
所有选择都必须通过“原始”字体特性选择完成。示例如下：
\begin{Verbatim}
  \fontspec{Charis SIL}[
    Renderer=Graphite,
    RawFeature={Uppercase Eng alternates=Large eng on baseline}]
  Ŋ
\end{Verbatim}

另一个示例：
\begin{Verbatim}
\fontspec{AwamiNastaliq-Regular.ttf}[Renderer=Graphite] ^^^^06b5
\addfontfeature{RawFeature={Lam with V=V over bowl}}    ^^^^06b5
\end{Verbatim}



\section{\MacOSX\ 的 \AAT\ 字体}
\label{sec:aat-features}

\begin{quote}\itshape
\textbf{警告！}
\XeTeX\ 在 \MacOSX\ 上的实现目前仍处于变动之中，下文所载信息很可能自
2013 年起就已过时。随着 \XeTeX\ 完成向纯跨平台应用的过渡，
本节所述特性很可能不再可用。本节中的所有示例现已移除。
\end{quote}

\MacOSX\ 的字体技术诞生于 OpenType 普及之前，围绕苹果发明的“\AAT”字体格式展开。
该格式有一些优点（也有其他缺点），但从未在字体界广泛流行。

尽管如此，这是 \XeTeX\ 最早支持的字体格式（究其根源，它最初就诞生于
\MacOSX），也是 \pkg{fontspec} 支持的第一种字体格式。随 \MacOSX\ 分发的
一些字体仍为 \AAT\ 格式，例如“Skia”。

\subsection{连字}

\feat{Ligatures} 指出于功能或美学原因，用专门绘制的字形替换两个独立字符。
对于 \AAT\ 字体，可从 \opt{Required}、\opt{Common}、\opt{Rare}（或
\opt{Discretionary}）、\opt{Logos}、\opt{Rebus}、\opt{Diphthong}、\opt{Squared}、
\opt{AbbrevSquared} 与 \opt{Icelandic} 的任意组合中选择。

其他一些苹果 \AAT\ 字体把那些“Rare”连字包含在 \opt{Icelandic} 特性中。
另请注意，在 \XeTeX\ 中，用空花括号对拆分连字这一旧式 \TeX\ 技巧并不奏效；
若不希望执行连字，必须使用 0\,pt 字距或 \cs{hbox}（\eg\ \cs{null}）来拆分字符
（常见适用情形是“shelf\null full”这类词）。

\subsection{字母} \label{sec:aat-letters}
\opt{Letters} 特性指定当前字体中字母的外观。对于 \AAT\ 字体，
可选择 \opt{Normal}、\opt{Uppercase}、\opt{Lowercase}、\opt{SmallCaps} 与
\opt{InitialCaps}。


\subsection{数字}
\feat{Numbers} 特性定义所选字体中数字的外观。对于 \AAT\ 字体，
可为 \opt{Lining} 或 \opt{OldStyle} 与 \opt{Proportional} 或
\opt{Monospaced} 的组合（后者适合表格材料）。同义词 \opt{Uppercase} 与
\opt{Lowercase} 分别等价于 \opt{Lining} 与 \opt{OldStyle}。
差异已在 \vref{sec:addfontfeatures} 中展示。

\subsection{上下文相关} \label{sec:contextuals}
此特性指随位置变化的字形替换；上下文花体等即在此实现。
\AAT\ 字体的选项有 \opt{WordInitial}、\opt{WordFinal}、\opt{LineInitial}、
\opt{LineFinal} 与 \opt{Inner}（有时称为“非末尾”）。与连字一样，
它们是非排他选择器，可在名称前加 \opt{No} 前缀将其关闭。

\subsection{上下位置}
\feat{VerticalPosition} 特性用于访问下标（\opt{Inferior}）、上标
（\opt{Superior}）数字与字母（有时还有少量标点）。\opt{Ordinal} 选项
（理应）具备上下文敏感性，仅提升紧跟数字之后出现的字符。

\pkg{realscripts} 宏包重定义了 \cmd\textsubscript\ 与 \cmd\textsuperscript\ 命令，
使其使用上述字体特性，包括在脚注标签中使用。

\subsection{分数}
许多字体具备排版各种形式分数材料的能力。在 \pkg{fontspec} 中通过
\feat{Fractions} 特性访问，在 \AAT\ 与 OpenType 字体中均可将其设为
\opt{On} 或 \opt{Off}。

在 \AAT\ 字体中，应使用“分数斜杠”（solidus）字符来创建分数。
当 \feat{Fractions} 设为 \opt{On} 时，仅使用预先绘制的分数。

使用 \opt{Diagonal} 选项（仅 \AAT）时，字体会尝试用上标与下标字符拼出分数。

另有一些字体（主要是亚洲字体）还提供 \opt{Alternate} 特性。

\subsection{变体}
\feat{Variant} 特性接受单个数值输入，用于选择不同字母字形。
\vref{sec:newfeatures} 介绍了为变体命名的方式，这应按字体逐一进行。

\subsection{替代字形}
\feat{Alternate} 的选择同样必须以数值进行。
\vref{sec:newfeatures} 介绍了为替代字形命名的方式，这应按字体逐一进行。


\subsection{风格}
\feat{Style} 特性的选项在 \AAT\ 中定义为下列之一：\opt{Display}、
\opt{Engraved}、\opt{IlluminatedCaps}、\opt{Italic}、
\opt{Ruby}\footnotemark、\opt{TallCaps} 或 \opt{Titling}。
\footnotetext{“Ruby”指小型光学字号，在日文排版中用于标注。}

这些特性的典型示例见 \ref{sec:ot-feat-style}。

\subsection{中日韩字形}
关于中日韩表意文字“应当”长什么样，历来有多种标准。一些字体包含大量替代字形，
以便在合适的任一种形式下正确显示这些字符。\AAT\ 与 OpenType 字体都支持下列
\feat{CJKShape} 选项：\opt{Traditional}、\opt{Simplified}、\opt{JIS1978}、
\opt{JIS1983}、\opt{JIS1990} 与 \opt{Expert}。OpenType 还支持 \opt{NLC} 选项。

\subsection{字符宽度}
相关示例见 \vref{sec:CharacterWidth}；OpenType 与 \AAT\ 字体的这些特性相同。
\AAT\ 还允许 \feat{CharacterWidth}|=|\opt{Default} 以返回原始字体设置。

\subsection{附加符号}
附加符号是加在字母上的标记，如锐音符或波浪号；通常表示读音变化。
在阿拉伯文字母中，附加符号用于表示元音。在 \AAT\ 字体中，
可选择 \opt{Show}、\opt{Hide} 或 \opt{Decompose}。
\opt{Hide} 选项用于阿拉伯文这类可带可不带元音符号的文字。\eg\
\verb|\fontspec[Diacritics=Hide]{...}|

随 \MacOSX\ 分发的一些较老字体，在 \feat{Diacritics} 特性标签下把“|O/|”
\etc\ 作为书写“\O”的简写。若遇到这类字体，需要关闭此特性
（试想输入 |hello/goodbye| 却得到“hell\o goodbye”！）：
把附加符号中的两个字符分解为实际想要的字符。我建议改用规范的
\LaTeX\ 输入约定来获取这类字符。


\subsection{圈码}
各种亚洲字体配备了更丰富的不同形式数字与数词。这些通过 \feat{Annotation}
特性访问，选项如下：\opt{Off}、\opt{Box}、\opt{RoundedBox}、\opt{Circle}、
\opt{BlackCircle}、\opt{Parenthesis}、\opt{Period}、\opt{RomanNumerals}、
\opt{Diamond}、\opt{BlackSquare}、\opt{BlackRoundSquare} 与 \opt{DoubleCircle}。

\end{document}

% /©
% ------------------------------------------------
% The FONTSPEC package  <latex3.github.io/fontspec>
% ------------------------------------------------
% Copyright  2022-2026  The LaTeX project,  LPPL "maintainer"
% Copyright  2004-2022  Will Robertson
% Copyright  2009-2015  Khaled Hosny
% Copyright  2013       Philipp Gesang
% Copyright  2013-2016  Joseph Wright
% ------------------------------------------------
% This package is free software and may be redistributed and/or modified under
% the conditions of the LaTeX Project Public License, version 1.3c or higher
% (your choice): <http://www.latex-project.org/lppl/>.
% ------------------------------------------------
% ©/
